\subsection{四个问题的共同结构}

四个问题看似分属检测、调度与装填三类，但它们共享同一个底层对象：\textbf{时频单元格}。
把计划 $P$ 展开为它所占用的单元集合
\[
  U(P)=\bigl\{(b,t)\ :\ f\le b<f+w,\ \ \exists\,k\in\{0,\dots,n-1\},\ s+k(d+g)\le t<s+d+k(d+g)\bigr\},
\]
则“两计划冲突”就等价于 $U(P_i)\cap U(P_j)\neq\varnothing$，“一组计划无冲突”等价于它们的单元集合两两不交。
这一观察把四个问题统一成同一语言：

\begin{itemize}[leftmargin=1.6em,itemsep=0.15em]
  \item \textbf{问题 1} 是判定 $\binom{\valQonePlans}{2}$ 对单元集合是否相交。直接求交代价高，但由于 $U(P)$ 是
        “频段区间 $\times$ 周期性时间区间”的乘积结构，相交性可用区间算术在常数时间内判定（见命题~\ref{prop:cell}）。
  \item \textbf{问题 2、4} 是：每个计划从一个有限动作集合中\emph{恰选一个}动作（保持、频移、时移、\,问题 4 另加改间隔、撤销），
        使得所选动作对应的单元集合两两不交，并在题目给出的优先级下最优。
        这是一个\emph{列表着色 / 带选项的装填}型问题，其特例（只许保持或撤销）正是冲突图的最大独立集，
        故为 NP-难\cite{karp1972,garey1979}。本题属于频率指配问题族\cite{hale1980,aardal2007}，
        但与经典频率指配（给每个发射机指定频点、相邻发射机须相隔足够频距）相比多出\emph{时间}这一维，
        且调整方式受“只改一个参数”的强限制。
  \item \textbf{问题 3} 是：在问题 2 方案留下的自由单元中，放入尽可能多的同形状“砖块”（C 类计划的单元集合），
        这是二维\emph{集合装填}（set packing）问题\cite{padberg1973,lodi2002}。
\end{itemize}

统一的语言带来一个直接的好处：无论求解器给出什么，都可以用同一个\emph{与求解器代码完全分离}的校验器
——把每个计划重新展开成单元集合、用集合求交重新找冲突——来复核。本文的全部结论都经过这一关。

\subsection{建模的三个关键抉择}

\paragraph{（一）冲突的精确语义。}题目只说“时间区间和频段区间都存在交叠”，未言明区间端点开闭、
是否逐次比较。附件的示例（$[35,40)$、间隔 $60$、第二次 $[100,105)$）给出了半开区间与推进规则，
本文据此把冲突定义在\emph{整个使用序列}上（MDR-0001）：只要存在一对使用次在时间上交叠且频段交叠，两计划即冲突。
若改用闭区间，相邻不重叠的两次使用会被误判为冲突；若只比较首次使用，“间隔时长”与“使用次数”两个参数将失去意义。

\paragraph{（二）优先目标如何形式化。}题目按“撤销 $\succ$ 调整数量 $\succ$ 类别优先级 $\succ$ 调整幅度”的次序给出四条目标，
却没有给出权重。把“类别优先级”嵌进前两条，得到本文的主方案 P：七级词典序
\[
  L=(\text{撤销}_A,\ \text{撤销}_B,\ \text{撤销}_C,\ \text{调整}_A,\ \text{调整}_B,\ \text{调整}_C,\ \text{归一化幅度}),
\]
逐级最小化。这不是唯一的读法，因此本文同时求解另外三种方案（总量优先 T、不含优先级 S、纯加权 W）并给出对照，
用数据回答“优先级条款到底起了多大作用”（第~\ref{sec:q2schemes} 节）。这是一个\emph{建模决策}而非计算细节，
故全部备选与理由记录在 MDR-0003 中。

\paragraph{（三）问题 3 中“不限制平移幅度”指谁。}“平移”一词只对已有位置的计划有意义，而新增装备本无“原位置”可平移，
故该句存在两种读法：解释 A（新增计划的放置不受 $10\Delta f/5\Delta t$ 约束，既有计划固定）与
解释 B（既有计划也可不受幅度限制地再次平移以腾出空间）。两种读法给出的答案相差约五倍。
本文以 A 为主解释（它与 \nolinkurl{result3.xlsx} 模板的三列一一对应，评阅者可仅凭结果文件独立复核），
同时为 B 给出可核查的上下界（第~\ref{sec:q3alt} 节、MDR-0010）。

\subsection{技术路线}

图~\ref{fig:roadmap} 给出本文的技术路线。论文中报告的全部数值都产生于同一次云端运行，
每个阶段的输入输出、参数、代码版本与 SHA-256 散列都写入清单文件，数字由该次运行的阶段产物自动注入（附录~\ref{app:repro}）。
唯一的例外是问题 2 与问题 4 求解时使用的\emph{现任解}（热启动，第~\ref{sec:q2alg} 节）：
它来自更早的一次运行，作为文件存放在仓库中，其来源与目标向量记录在文件内。本文对这一依赖如实报告，
并另跑一次不加载现任解的冷启动作对照。

\begin{figure}[htbp]
\centering
\begin{tikzpicture}[
  every node/.style={font=\small},
  box/.style={draw=ink!80,rounded corners=2pt,align=center,inner sep=4pt,
              minimum height=1.35cm,text width=3.2cm},
  data/.style={box,fill=ink!8},
  model/.style={box,fill=accent!12},
  check/.style={box,fill=black!6},
  arr/.style={-{Stealth[length=2.2mm]},thick,draw=ink!75}
]
% ---- row 1: detection
\node[data]  (r1a) at (0,0)    {附件 1\\$\valQonePlans$ 个用频计划};
\node[data]  (r1b) at (4.1,0)  {契约校验与形式化\\$P_i=(f_i,w_i,s_i,d_i,g_i,n_i)$};
\node[model] (r1c) at (8.2,0)  {\textbf{问题 1}\\成对枚举 $\;\|\;$ 分桶扫描线};
\node[model] (r1d) at (12.3,0) {冲突图 $G$\\度分布、分量、类别对};
\draw[arr] (r1a)--(r1b); \draw[arr] (r1b)--(r1c); \draw[arr] (r1c)--(r1d);
% ---- row 2: resolution
\node[model] (r2a) at (0,-2.9)    {动作集合 $O_i$\\单元格 AtMostOne 模型};
\node[model] (r2b) at (4.1,-2.9)  {\textbf{问题 2}\\七级词典序 CP-SAT\\（现任解上界割）};
\node[model] (r2c) at (8.2,-2.9)  {\textbf{问题 4}\\C 类增加改间隔动作};
\node[data]  (r2d) at (12.3,-2.9) {无冲突方案\\$\mathcal{S}_2$、$\mathcal{S}_4$};
\draw[arr] (r2a)--(r2b); \draw[arr] (r2b)--(r2c); \draw[arr] (r2c)--(r2d);
\draw[arr] (r1d.south) -- ++(0,-0.6) -| (r2a.north);
% ---- row 3: packing, verification, deliverables
\node[model] (r3a) at (0,-5.8)    {\textbf{问题 3}\\自由单元 $\to$ 候选位置\\集合装填 CP-SAT};
\node[model] (r3b) at (4.1,-5.8)  {上界链\\CP-SAT / LP / 面积守恒};
\node[check] (r3c) at (8.2,-5.8)  {独立校验器\\单元集合求交、\\穷举对照、灵敏度};
\node[data]  (r3d) at (12.3,-5.8) {\texttt{result1--4.xlsx}\\论文图、表、数字};
\draw[arr] (r3a)--(r3b); \draw[arr] (r3b)--(r3c); \draw[arr] (r3c)--(r3d);
\draw[arr] (r2d.south) -- ++(0,-0.6) -| (r3a.north);
% ---- verification feeds back on every model block
\draw[arr,dashed,draw=black!55] (r1c.south) -- ++(0,-0.45) -| ([xshift=-1.1cm]r3c.north);
\draw[arr,dashed,draw=black!55] (r2c.south) -- ([xshift=0cm]r3c.north);
\end{tikzpicture}
\caption{技术路线。实线为数据流，虚线为独立校验（校验器不复用检测与求解代码，
把每个计划重新展开为时频单元集合后用集合求交复核）。灰蓝色块为数据与产物，橙色块为模型与算法。}
\label{fig:roadmap}
\end{figure}
